\label{sec:immunity}

We prove that all four divisor methods --- Huntington-Hill, Webster,
Adams, and Jefferson --- are immune to all three classical apportionment
paradoxes.
The proof is unified: it applies to any divisor method.

\subsection{Unified Divisor Method Framework}

A divisor method is characterised by a rounding function $r(q)$
mapping each exact quota $q$ to an integer seat count.
The rounding function is non-decreasing and satisfies $r(q) \in
\{\lfloor q \rfloor, \lceil q \rceil\}$ for most methods
(HH, Webster, Adams, Jefferson each satisfy this for integer-valued
divisors).

Equivalently, a divisor method uses a priority function $P_i(s)$
that is strictly decreasing in $s$ for fixed $p_i$
and strictly increasing in $p_i$ for fixed $s$.
All four divisor methods satisfy:
\begin{itemize}
  \item $P_i(s)$ is strictly decreasing in $s$ (each additional seat
    has lower priority than the previous).
  \item $P_i(s)$ is strictly increasing in $p_i$ (larger population
    gives higher priority for each seat).
  \item $P_i(s) \to 0$ as $s \to \infty$ (no state takes all seats).
  \item $P_i(0) \to \infty$ or $P_i(0) = p_i$ (all states get at least 1 seat
    before any state gets its second seat).
\end{itemize}

\subsection{Theorem: Divisor Methods Are Paradox-Free}

\begin{theorem}[Divisor Method Paradox Immunity]
  \label{thm:divisor-immune}
  Any divisor method with priority function $P_i(s)$ that is
  (a)~strictly decreasing in $s$ and (b)~non-decreasing in $p_i$
  satisfies:
  \begin{enumerate}
    \item[(i)] \textbf{Alabama paradox immunity}: If $H$ increases to $H+1$,
      no state's seat count decreases.
    \item[(ii)] \textbf{Population paradox immunity}: If $p_i$ increases
      (all other populations fixed, $H$ fixed), state $i$'s seat count
      does not decrease.
    \item[(iii)] \textbf{New-states paradox immunity}: Adding state $k$
      with $m$ new seats (where $m \approx p_k \times H / N$,
      $H$ increases to $H+m$) does not change the seat count of
      any existing state.
  \end{enumerate}
\end{theorem}

\begin{proof}
  \textbf{(i) Alabama paradox immunity.}
  The priority-queue algorithm for $H$ seats assigns seat $t$ to the
  state with highest priority at step $t$.
  The allocation for seats $1, 2, \ldots, H$ is fixed: it depends
  only on the population vector $\vec{p}$ and the priority function.
  When $H$ increases to $H+1$, seat $H+1$ is assigned to the state
  with highest priority after the first $H$ seats have been allocated.
  This operation adds exactly one seat to exactly one state and
  leaves all other seat counts unchanged.
  In particular, no state's seat count decreases.

  \textbf{(ii) Population paradox immunity.}
  Let $\vec{p}' = \vec{p} + \delta_i \mathbf{e}_i$ where $\delta_i > 0$
  (state $i$'s population increases by $\delta_i$).
  Under $\vec{p}'$, state $i$'s priority $P_i(s)$ increases for all $s$
  (by assumption (b)).
  All other states' priorities are unchanged.
  The priority sequence for allocating $H$ seats is thus at least as
  favourable to state $i$ under $\vec{p}'$ as under $\vec{p}$:
  wherever state $i$ appeared in the priority sequence under $\vec{p}$,
  it appears at the same position or earlier under $\vec{p}'$.
  Therefore, state $i$ receives at least as many seats under $\vec{p}'$
  as under $\vec{p}$.

  More precisely, if state $i$ received $s_i$ seats under $\vec{p}$,
  then under $\vec{p}'$, state $i$'s priority for its $(s_i+1)$-th seat
  may now exceed some other state's priority, giving it $s_i+1$ or more
  seats.
  But it cannot lose seats, because its priorities for seats
  $1, \ldots, s_i$ all increased, so all these seats are still
  awarded to state $i$ before any competing state.

  \textbf{(iii) New-states paradox immunity.}
  Adding state $k$ with population $p_k$ to an $n$-state allocation
  with House size $H$, and increasing $H$ to $H + m$ where
  $m$ is chosen so that state $k$ receives exactly $m$ seats
  in the priority queue order:

  The priority sequence for the original $n$ states and $H$ seats
  is a well-defined sequence of (state, seat-number) pairs.
  Inserting state $k$ with $m$ new seats means adding $m$ new entries
  to this sequence at the positions determined by $k$'s priority values.
  The key claim is that these insertions do not perturb the relative
  ordering of the original $n$ states' priorities.

  Formally, for any two original states $i \neq j$ and seat numbers
  $s_i, s_j$, the comparison $P_i(s_i) > P_j(s_j)$ depends only on
  $p_i, p_j, s_i, s_j$ and is unaffected by $p_k$.
  The only change is that some of the ``slots'' in the priority
  sequence for original states' seats are occupied by state $k$'s
  seats --- but this can only shift original states' seats forward
  in the sequence (making them arrive at a higher slot number
  rather than a lower one), not change which seats they receive
  relative to each other.

  Therefore, the seat allocation among original states $1, \ldots, n$
  under $H + m$ seats is the same as under $H$ seats (with $m$
  additional seats going to state $k$).
  No existing state's seat count changes.

  A subtlety: the divisor $d$ adjusts to maintain the total at $H+m$,
  and this adjustment can in principle shift some states' allocations.
  The formal argument requires showing that the divisor adjustment
  is small enough not to perturb any rounding threshold.
  This is guaranteed when $m/H \approx p_k/N$ (state $k$ gets its
  exact proportional share of seats), which is the definition of
  a ``proportional'' new-state addition.
  See \citet{balinski1982fair}, ch.~5, for the technical details.
\end{proof}

\subsection{Verification: Huntington-Hill on Census Data}

The \textsc{bisect-apportion} crate empirically checks Alabama-paradox immunity
for all four implemented divisor methods on synthetic population fixtures:
\begin{itemize}
  \item \textbf{Alabama paradox check}: For house sizes 5 through 50 on a
    five-state synthetic population vector, no state's seat count decreases as
    $H$ increases under Huntington-Hill, Webster, Adams, or Jefferson.
  \item \textbf{House-size monotonicity}: A three-state synthetic fixture checks
    that each state's seat series is non-decreasing across house sizes 3 through
    30 for all four divisor rules.
\end{itemize}
These checks support the formal proof but do not replace future Census-year
fixtures.

\subsection{Hash-Bound Divisor Smoke Package}

The repository now includes a minimal positive divisor-method evidence package at
\texttt{docs/examples/j-apportionment-evidence-packages/divisor-method-smoke}.
Its manifest uses schema \texttt{j-divisor-evidence-manifest v1}; it hash-binds
the three-state allocation fixture and replays Webster, Adams, and Jefferson
through \texttt{bisect-apportion::divisor\_evidence}.  The same fixture also
checks the shared divisor Alabama-paradox surface across Huntington-Hill,
Webster, Adams, and Jefferson for house sizes 3 through 10, expecting no
decreasing seat events.
